%
% File: introduction.tex
%
\chapter{Introduction}
\label{chap:intro}
\setlength{\parskip}{1em}
%
The classroom of 2026 is multilingual, neurodiverse, and increasingly mediated by artificial intelligence. In a single lecture hall, an instructor may simultaneously address learners who are still acquiring the language of instruction, learners with diagnosed conditions such as dyslexia or attention-deficit/hyperactivity disorder, and learners whose only difficulty is the unevenness of prior schooling. For each of these students, a textbook paragraph or a verbal explanation imposes a cognitive cost that is not necessarily related to the material itself but to the linguistic packaging of that material. The pedagogical problem is therefore not what is taught, but how it is encoded for transmission. Generative artificial intelligence, and especially the recent generation of multimodal large language models capable of producing diagrams, illustrations, and short animations from natural-language prompts, offers a route around this bottleneck. Rather than asking learners to convert text into mental images on their own --- a task that is itself language-dependent --- the technology can supply the visual representation directly, at the right level of abstraction, and within seconds.

The emergence of tools such as GPT-4 Vision, DALL$\cdot$E~3, and Gemini has shifted the conversation about AI in education away from text-only chatbots and towards multimodal systems that read, write, see, and draw. \citet{bewersdorff2024} characterise this as a ``next step'' in generative AI, in which the same model can interpret a hand-drawn student diagram and respond with both a corrected diagram and a level-appropriate textual explanation. Earlier waves of educational technology --- intelligent tutoring systems, adaptive testing platforms, screen readers --- addressed accessibility one channel at a time. Multimodal AI collapses these channels into a single interface and, crucially, makes visual content production cheap enough to be personalised. A teacher who would once have spent hours preparing differentiated worksheets can now generate a learner-specific diagram during the lesson itself. The pedagogical implications are significant for any student who benefits from non-verbal scaffolding, but they are particularly consequential for learners with low language proficiency and for learners with disabilities, two groups for whom the gap between cognitive capability and linguistic access is often the primary barrier to achievement.

The case for visual explanations is not new. Mayer's Cognitive Theory of Multimedia Learning (CTML) and Sweller's Cognitive Load Theory (CLT), both of which have been refined for over three decades, establish that learners process information through partially independent verbal and visual channels and that learning is optimised when extraneous load is minimised and germane load is maximised. What is new is the scale and immediacy with which AI can deliver visual material that respects these principles. \citet{twabu2025} argue that traditional CTML and CLT, although foundational, were articulated in an era of static multimedia and require extension to accommodate the adaptive, learner-responsive behaviour of AI systems. Their proposed AI-augmented framework is one of several recent attempts to bring multimedia learning theory up to date with the technological reality of generative models.

Despite a rapidly expanding literature on AI in education, work that specifically interrogates the pedagogical potential of AI-\emph{generated} visual explanations for learners with low language proficiency or learning disabilities remains scattered across systematic reviews of special education technology, theoretical papers on multimodal AI, and empirical studies of individual tools. No synthesis to date has brought these strands together under a single analytic frame. The present review addresses that gap by examining ten studies published between 2024 and 2025 in light of the following research questions. \textbf{RQ1.} What categories of AI-generated visual explanations have been documented in recent pedagogical research relevant to learners with low language proficiency or learning disabilities? \textbf{RQ2.} How do these tools and approaches align with the principles of CTML and CLT? \textbf{RQ3.} What evidence exists for their pedagogical efficacy with the target learner populations? \textbf{RQ4.} What limitations and unresolved questions does the current literature reveal, and what should future research prioritise?

The significance of this review is threefold. First, it consolidates a fragmented evidence base at a moment when generative AI is being adopted in classrooms faster than the research can evaluate. Second, it offers educators, instructional designers, and policymakers a theoretically grounded basis for decisions about which AI tools to introduce, for whom, and under what conditions. Third, by foregrounding learners with low language proficiency alongside learners with disabilities, it resists the tendency of the field to treat these populations as separate, when in practice they share the same underlying need for non-verbal scaffolding. The remainder of the paper is organised as follows. Chapter~\ref{chap:litreview} reviews the relevant literature on multimedia learning theory, AI applications for diverse learners, and visual generation in educational contexts. Chapter~\ref{chap:method} sets out the systematic review methodology. Chapter~\ref{chap:findings} presents the thematic findings. Chapter~\ref{chap:discussion} discusses theoretical and practical implications, acknowledges limitations, and proposes directions for future inquiry.
